\documentclass[12pt,a4paper]{article}

\usepackage{iftex}
\ifPDFTeX
  \usepackage[T2A]{fontenc}
  \usepackage[utf8]{inputenc}
  \usepackage[russian]{babel}
\else
  \usepackage{fontspec}
  \usepackage{polyglossia}
  \setmainlanguage{russian}
  \setmainfont{DejaVu Serif}
  \setsansfont{DejaVu Sans}
  \setmonofont{DejaVu Sans Mono}
\fi
\usepackage{amsmath,amssymb,amsthm,bm,mathtools}
\usepackage{geometry}
\usepackage{enumitem}
\usepackage{hyperref}
\usepackage{physics}
\usepackage{xcolor}
\geometry{margin=2.2cm}
\hypersetup{colorlinks=true,linkcolor=blue,citecolor=blue,urlcolor=blue}
\setlist[itemize]{leftmargin=1.5em}
\setlist[enumerate]{leftmargin=1.7em}
\setlength{\parindent}{1.25em}
\setlength{\parskip}{0.35em}

\newcommand{\E}{\mathcal{E}}
\newcommand{\Hatom}{\mathcal{H}_{\mathrm{at}}}
\newcommand{\Hsp}{\mathcal{H}_{\mathrm{sp}}}
\newcommand{\m}{\bm{m}}
\newcommand{\ez}{\bm{e}_z}
\newcommand{\ta}{\tau_{\mathrm{a}}}
\newcommand{\tm}{\tau_{\mathrm{m}}}
\newcommand{\R}{\bm{R}}
\newcommand{\uvec}{\bm{u}}
\newcommand{\rvec}{\bm{r}}
\newcommand{\ezz}{u_{zz}}

\begin{document}

\begin{center}
{\LARGE Магнитный бризер, запускаемый атомным бризером:\\[0.2em]
редуцированная магнитоупругая постановка в физических переменных}\\[0.8em]
{\large Рабочая заметка на основе исходной математической схемы и статьи по магнитоупругому учету деформации}\\[0.6em]
\end{center}

\section*{Аннотация}
Ниже переписана исходная постановка задачи так, чтобы магнитоупругая связь была введена в более физической форме: не через абстрактную модуляцию коэффициентов ``от смещений'', а через изменение межузельных расстояний и соответствующие производные обменных и DM-констант по деформации. Для одноионной анизотропии используется не разложение $A(\eta)$, а стандартный линейный магнитоупругий вклад с постоянной константой связи $B_A$, так что временная зависимость возникает только через локальную координату деформации. По сути задача остается той же: сначала независимо рассчитывается атомный бризер в одномерной решетке, затем его поле смещений используется как внешний локализованный драйв для спиновой подсистемы. Изменяется именно язык записи. Вместо чисто математических разностей $q_{n+1}-q_n$ вводятся равновесные положения узлов $R_n$, текущие положения $r_n$, деформированные длины связей $\ell_n$ и относительная деформация $\eta_n=(\ell_n-a)/a$, а магнитоупругая часть строится по тому же принципу, что и в работе о распространении ультразвука и exchange-strain механизме \cite{Tereshchenko2022}. Это позволяет прямо видеть, какие именно производные $\dv*{J}{\ell}$ и $\dv*{D}{\ell}$ входят в модель, а анизотропный канал выделяется через отдельную магнитоупругую константу $B_A$.

\section{Идея редукции и физические переменные}
Рассматривается одномерная атомная цепочка с равновесными положениями
\begin{equation}
R_n = na, \qquad n\in\mathbb{Z},
\end{equation}
где $a$ --- шаг решетки вдоль оси цепочки $z$. Текущее положение $n$-го узла записываем как
\begin{equation}
r_n(t)=R_n+u_n(t),
\end{equation}
где $u_n(t)$ --- продольное смещение узла.

После вычисления атомного бризера $u_n^{(b)}(t)$ мы не решаем полную самосогласованную магнитоупругую задачу, а используем этот профиль как заданное поле деформации, которое модулирует параметры спинового гамильтониана. Таким образом:
\begin{align}
&\text{(I) атомная задача:} && u_n^{(b)}(t) \ \text{--- периодический локализованный бризер}, \\
&\text{(II) магнитная задача:} && \m_n(t) \ \text{--- спиновая динамика в неавтономных } J_n(t),D_n(t),A_n(t).
\end{align}

Ключевой принцип тот же, что и в исходной заметке, но теперь драйв формируется через \emph{изменение расстояния между узлами}, а не через формально введенную модуляцию коэффициентов.

\section{Атомная подсистема и деформированная геометрия}
В одномерной FPU-$\beta$ цепочке удобно писать гамильтониан через смещения $u_n$:
\begin{equation}
\Hatom = \sum_n \frac{M\dot u_n^2}{2} + \sum_n V\bigl(r_{n+1}-r_n\bigr),
\qquad
V(\Delta r)=\frac{K}{2}(\Delta r-a)^2+\frac{\beta}{4}(\Delta r-a)^4.
\label{eq:Hatom_physical}
\end{equation}
Так как
\begin{equation}
r_{n+1}-r_n = a + u_{n+1}-u_n,
\end{equation}
то в безразмерной записи это эквивалентно использованной ранее модели через разности $q_{n+1}-q_n$.

Физически важная величина --- это не само смещение узла, а деформированная длина связи
\begin{equation}
\ell_n(t)=r_{n+1}(t)-r_n(t)=a+u_{n+1}(t)-u_n(t)
\label{eq:ell_n}
\end{equation}
и относительная деформация связи
\begin{equation}
\eta_n(t)=\frac{\ell_n(t)-a}{a}=\frac{u_{n+1}(t)-u_n(t)}{a}.
\label{eq:eta_n_discrete}
\end{equation}
В длинноволновом или квазиконтинуальном языке это просто одномерная компонента тензора деформации,
\begin{equation}
\eta_n(t)\approx \eval{\pdv{u(z,t)}{z}}_{z=R_n+a/2},
\qquad
\ezz(z,t)=\pdv{u(z,t)}{z},
\label{eq:strain_continuum}
\end{equation}
что полностью аналогично записи через производные смещений, используемой в \cite{Tereshchenko2022}.

Для site-параметров естественно использовать деформацию в узле,
\begin{equation}
\bar\eta_n(t)=\frac{\eta_n(t)+\eta_{n-1}(t)}{2}
=\frac{u_{n+1}(t)-u_{n-1}(t)}{2a}
\approx \eval{\pdv{u(z,t)}{z}}_{z=R_n}.
\label{eq:eta_site}
\end{equation}

Если атомная подсистема имеет бризерное решение $u_n^{(b)}(\ta)$ периода $T_{\mathrm a}$, то индуцированное поле деформации определяется как
\begin{equation}
\eta_n^{(b)}(\ta)=\frac{u_{n+1}^{(b)}(\ta)-u_n^{(b)}(\ta)}{a},
\qquad
\bar\eta_n^{(b)}(\ta)=\frac{u_{n+1}^{(b)}(\ta)-u_{n-1}^{(b)}(\ta)}{2a}.
\label{eq:breather_strain}
\end{equation}
Разложим их на статическую и осциллирующую части:
\begin{equation}
\eta_n^{(b)}(\ta)=\eta_n^{(0)}+\delta\eta_n(\ta),
\qquad
\bar\eta_n^{(b)}(\ta)=\bar\eta_n^{(0)}+\delta\bar\eta_n(\ta),
\label{eq:split_eta}
\end{equation}
где средние по периоду величины отвечают за статическую перенормировку магнитных параметров, а переменные части играют роль драйва.

\section{Связь атомного смещения с магнитными константами}
В статье \cite{Tereshchenko2022} exchange-strain взаимодействие вводится через зависимость обмена от расстояния между магнитными центрами:
\begin{equation}
J\bigl(R_{ij}+u_i-u_j\bigr)-J(R_{ij})
\simeq
\eval{\dv{J}{\ell}}_{\ell=R_{ij}} (u_i-u_j).
\label{eq:exchange_article_style}
\end{equation}
Для нашей одномерной цепочки тот же принцип дает
\begin{equation}
J_n(t)=J\bigl(\ell_n(t)\bigr),
\qquad
D_n(t)=D\bigl(\ell_n(t)\bigr),
\qquad
A_n(t)=A_0+B_A\,\bar\eta_n(t).
\label{eq:JDA_functional}
\end{equation}

То есть обмен и DM-константа рассматриваются как функции деформированной длины связи, а одноионная анизотропия получает обычную линейную магнитоупругую поправку с постоянной константой $B_A$, не зависящей от времени. Временная зависимость $A_n(t)$ возникает только потому, что сама локальная деформация $\bar\eta_n(t)$ создается движением узлов.

Разложение по малой деформации около недеформированного состояния тогда принимает вид
\begin{align}
J_n(t) &\simeq J_0 + \eval{\dv{J}{\ell}}_{a}\,\bigl(\ell_n(t)-a\bigr)
= J_0 + a\,J_\ell'\,\eta_n(t),
\label{eq:J_expand}\\
D_n(t) &\simeq D_0 + \eval{\dv{D}{\ell}}_{a}\,\bigl(\ell_n(t)-a\bigr)
= D_0 + a\,D_\ell'\,\eta_n(t),
\label{eq:D_expand}\\
A_n(t) &= A_0 + B_A\,\bar\eta_n(t).
\label{eq:A_expand}
\end{align}
Здесь
\begin{equation}
J_\ell'=\eval{\dv{J}{\ell}}_{\ell=a},
\qquad
D_\ell'=\eval{\dv{D}{\ell}}_{\ell=a},
\end{equation}
а $B_A$ --- постоянная магнитоупругая константа, связывающая одноионную анизотропию с локальной site-деформацией.

В безразмерной форме удобно ввести коэффициенты чувствительности для bond-каналов и нормированную магнитоупругую константу для анизотропии:
\begin{equation}
\gamma_J = \frac{a}{J_0}\eval{\dv{J}{\ell}}_{a},
\qquad
\gamma_D = \frac{a}{2J_0}\eval{\dv{D}{\ell}}_{a},
\qquad
\gamma_A = \frac{B_A}{2J_0}.
\label{eq:gamma_defs}
\end{equation}
Тогда вариации параметров непосредственно выражаются через бризерную деформацию.

\section{Согласование времен атомной и магнитной задач}
Как и в исходной заметке, сопоставлять нужно одно и то же физическое время $t$:
\begin{equation}
 t = \tau_0^{(\mathrm a)}\,\ta = t_0^{(\mathrm m)}\,\tm,
 \qquad
 t_0^{(\mathrm m)}=\frac{\hbar}{2J_0 S}.
\label{eq:physical_time_relation_new}
\end{equation}
Отсюда
\begin{equation}
\ta = \rho\,\tm,
\qquad
\rho = \frac{t_0^{(\mathrm m)}}{\tau_0^{(\mathrm a)}}.
\label{eq:rho_def_new}
\end{equation}
Если атомный бризер имеет частоту $\Omega_{\mathrm a}$ в своей переменной $\ta$, то в магнитных единицах драйв действует с частотой
\begin{equation}
\omega_{\mathrm d}=\rho\,\Omega_{\mathrm a}.
\label{eq:omega_drive}
\end{equation}
Именно $\rho$ остается главным параметром резонансного сканирования.

\section{Редуцированный спиновый гамильтониан в физической форме}
Для магнитной подсистемы оставляем ту же структуру гамильтониана, что и в заметке и в работе о магнитных dark breathers, но заменяем коэффициенты на деформационно-модулированные:
\begin{equation}
\Hsp(t)= -2\sum_n J_n(t)\,\bm S_n\cdot\bm S_{n+1}
+ \sum_n A_n(t)\,(S_n^z)^2
+ \sum_n D_n(t)\,[\bm S_n\times \bm S_{n+1}]_z
- H\sum_n S_n^z.
\label{eq:Hspin_physical}
\end{equation}

Подставляя \eqref{eq:J_expand}--\eqref{eq:A_expand}, получаем
\begin{align}
J_n(t) &= J_0\Bigl[1+\gamma_J\,\eta_n^{(b)}\bigl(\rho\tm\bigr)\Bigr],
\label{eq:J_final}\\
D_n(t) &= D_0 + 2J_0\gamma_D\,\eta_n^{(b)}\bigl(\rho\tm\bigr),
\label{eq:D_final}\\
A_n(t) &= A_0 + 2J_0\gamma_A\,\bar\eta_n^{(b)}\bigl(\rho\tm\bigr)
= A_0 + B_A\,\bar\eta_n^{(b)}\bigl(\rho\tm\bigr).
\label{eq:A_final}
\end{align}
Если требуется избежать искусственного начального удара, можно умножить все добавки на адиабатический множитель
\begin{equation}
\chi(\tm)=1-e^{-\tm/\tau_{\mathrm{on}}}.
\end{equation}
Тогда в \eqref{eq:J_final}--\eqref{eq:A_final} каждая поправка просто умножается на $\chi(\tm)$.

После перехода к безразмерным спинам $\m_n=\bm S_n/S$ и энергии $\E=\Hsp/(2J_0S^2)$ имеем
\begin{equation}
\E = -\sum_n j_n(\tm)\,\m_n\cdot\m_{n+1}
+ \sum_n b_n(\tm)\,(m_n^z)^2
+ \sum_n d_n(\tm)\,[\m_n\times\m_{n+1}]_z
- h\sum_n m_n^z,
\label{eq:E_dimensionless_new}
\end{equation}
где
\begin{align}
j_n(\tm) &= 1+\gamma_J\,\eta_n^{(b)}\bigl(\rho\tm\bigr), \\
d_n(\tm) &= d_0+\gamma_D\,\eta_n^{(b)}\bigl(\rho\tm\bigr),
\qquad d_0=\frac{D_0}{2J_0}, \\
b_n(\tm) &= b_0+\gamma_A\,\bar\eta_n^{(b)}\bigl(\rho\tm\bigr),
\qquad b_0=\frac{A_0}{2J_0}, \\
h&=\frac{H}{2J_0S}.
\end{align}

\section{Уравнения движения}
Для численного интегрирования удобно оставить векторную форму Ландау--Лифшица:
\begin{equation}
\dv{\m_n}{\tm} = -\m_n\times \bm h_n^{\mathrm{eff}}
-\alpha_G\,\m_n\times\bigl(\m_n\times \bm h_n^{\mathrm{eff}}\bigr),
\qquad
\bm h_n^{\mathrm{eff}}=-\pdv{\E}{\m_n}.
\label{eq:LL_equation_new}
\end{equation}
Из \eqref{eq:E_dimensionless_new} следует
\begin{equation}
\bm h_n^{\mathrm{eff}}=
 j_{n-1}\m_{n-1}+j_n\m_{n+1}
 + d_{n-1}(\ez\times \m_{n-1}) - d_n(\ez\times \m_{n+1})
 + \bigl(h-2b_n m_n^z\bigr)\ez.
\label{eq:heff_new}
\end{equation}
Структура уравнений не изменилась; изменилась только физическая интерпретация коэффициентов $j_n,d_n,b_n$.

\section{Минимальная модель и рекомендуемая полная модель}
Если хочется начать с самого простого теста, можно оставить только канал модуляции анизотропии:
\begin{equation}
j_n\equiv 1,
\qquad d_n\equiv d_0,
\qquad b_n(\tm)=b_0+\gamma_A\,\bar\eta_n^{(b)}(\rho\tm).
\label{eq:minimal_model_new}
\end{equation}
Это уже физичнее прежней записи $A_n\propto q_{n+1}-q_n$, потому что здесь анизотропия получает стандартную локальную магнитоупругую добавку, а ее временная зависимость целиком задается деформацией узла, выраженной через производную поля смещений.

Но основной рекомендуемой постановкой остается полная модель \eqref{eq:E_dimensionless_new}, поскольку:
\begin{enumerate}
    \item обмен и DM-взаимодействие зависят именно от межузельного расстояния и потому должны чувствовать bond-деформацию $\eta_n$;
    \item одноионная анизотропия получает локальную линейную магнитоупругую добавку $B_A\bar\eta_n$, где сама константа $B_A$ времени не зависит;
    \item атомный бризер создает локализованный профиль $\partial_z u$, а не глобальную однородную модуляцию, поэтому именно локальные $j_n,d_n,b_n$ соответствуют физике задачи.
\end{enumerate}

\section{Что именно изменилось по сравнению с исходной заметкой}
Содержательно почти ничего не меняется; переписана именно физическая оболочка формул.

\begin{enumerate}
    \item Вместо переменной $q_n$ теперь используются равновесные положения $R_n$, текущие положения $r_n=R_n+u_n$ и длины связей $\ell_n=r_{n+1}-r_n$.
    \item Драйв определяется через относительную деформацию $\eta_n=(\ell_n-a)/a\approx \partial_z u$, то есть через аналог деформационного тензора, как в статье \cite{Tereshchenko2022}.
    \item Магнитные коэффициенты вводятся как функции геометрии $J(\ell)$ и $D(\ell)$, а для одноионной анизотропии используется линейный магнитоупругий член $B_A\bar\eta_n$.
    \item Безразмерные коэффициенты связи приобретают прямой смысл: для обмена и DM это коэффициенты чувствительности, выраженные через производные по длине связи, а для анизотропии --- нормированная магнитоупругая константа.
\end{enumerate}

Именно в таком виде модель легче защищать как физическую: в ней явно видно, что атомный бризер сначала меняет локальную геометрию решетки, а уже эта геометрия меняет магнитный гамильтониан.

\section{Итоговая формулировка задачи для кода}
Для практического расчета удобно использовать следующую последовательность.
\begin{enumerate}
    \item Вычислить атомный бризер $u_n^{(b)}(\ta)$ в FPU-$\beta$ цепочке.
    \item По нему построить
    \begin{equation*}
    \eta_n^{(b)}(\ta)=\frac{u_{n+1}^{(b)}(\ta)-u_n^{(b)}(\ta)}{a},
    \qquad
    \bar\eta_n^{(b)}(\ta)=\frac{u_{n+1}^{(b)}(\ta)-u_{n-1}^{(b)}(\ta)}{2a}.
    \end{equation*}
    \item Использовать эти профили в \eqref{eq:J_final}--\eqref{eq:A_final} или в их безразмерной форме в \eqref{eq:E_dimensionless_new}.
    \item Интегрировать спиновую динамику по \eqref{eq:LL_equation_new} и искать захват в локализованный периодический режим.
\end{enumerate}

В таком варианте сохраняется вся логика исходной заметки, но запись становится существенно более физичной: вместо ``модуляции от разностей'' мы работаем с деформированной длиной связи, относительной деформацией, производными $J$ и $D$ по длине связи и прямым линейным магнитоупругим вкладом в анизотропию.

\begin{thebibliography}{9}
\bibitem{Tereshchenko2022}
A.~A.~Tereshchenko, A.~S.~Ovchinnikov, D.~I.~Gorbunov, D.~S.~Neznakhin,
\textit{Theory of ultrasound propagation in LuCo$_3$ near the low-spin--high-spin crossover},
Phys. Rev. B \textbf{106}, 054417 (2022).

\bibitem{Bostrem2021}
I.~G.~Bostrem, E.~G.~Ekomasov, J.~Kishine, A.~S.~Ovchinnikov, V.~E.~Sinitsyn,
\textit{Dark discrete breather modes in a monoaxial chiral helimagnet with easy-plane anisotropy},
Phys. Rev. B \textbf{104}, 214420 (2021).
\end{thebibliography}

\end{document}
